\documentclass[11pt]{article}
\usepackage{amsmath,amssymb}
\usepackage[a4paper,margin=1in]{geometry}
\usepackage{hyperref}

% Remove paragraph indentation
\setlength{\parindent}{0pt}
\setlength{\parskip}{6pt}

\begin{document}

\title{A Concise Derivation of t-SNE}
\author{Zhenwei Tang}
\date{}
\maketitle

% ============================================================
\section*{Goal}
% ============================================================

Given $n$ high-dimensional data points
\[
  x_i \in \mathbb{R}^p, \qquad i = 1, \dots, n,
\]
t-SNE (t-distributed Stochastic Neighbor Embedding) aims to find low-dimensional
embeddings
\[
  y_i \in \mathbb{R}^d, \qquad d \ll p \quad (\text{typically } d = 2 \text{ or } 3),
\]
such that \emph{neighborhood structure is preserved}: pairs of points that are
close in the original space should also be close in the embedding, and pairs that
are far apart should remain far apart.

The key idea is to turn pairwise distances into probability distributions both in
the high-dimensional space and in the low-dimensional space, then minimize the
discrepancy between these two distributions.

% ============================================================
\section*{Step 1: Define Pairwise Similarities in High-Dimensional Space}
% ============================================================

\textbf{Why:} Raw Euclidean distances are hard to compare across different local
densities. Converting them to conditional probabilities normalises for local
density and gives a clear probabilistic interpretation of ``neighbourhood''.

For each pair $(i, j)$, define the \emph{conditional probability} that point
$x_i$ would pick $x_j$ as its neighbour, modelled by a Gaussian centred at
$x_i$:
\[
  p_{j \mid i}
  \;=\;
  \frac{\exp\!\left(-\|x_i - x_j\|^2 \,/\, 2\sigma_i^2\right)}
       {\displaystyle\sum_{k \neq i} \exp\!\left(-\|x_i - x_k\|^2 \,/\, 2\sigma_i^2\right)},
  \qquad p_{i \mid i} = 0.
\]

Here $\sigma_i > 0$ is a per-point bandwidth, chosen automatically by the
\textbf{perplexity} criterion (see Appendix~A).

To obtain a \emph{symmetric} joint probability, average the two conditional
probabilities:
\[
  \boxed{p_{ij} \;=\; \frac{p_{j\mid i} + p_{i\mid j}}{2n}},
  \qquad \sum_{i \neq j} p_{ij} = 1.
\]

The symmetrisation ensures that $p_{ij} = p_{ji}$ and that every point has a
meaningful influence on every other point.

% ============================================================
\section*{Step 2: Define Pairwise Similarities in Low-Dimensional Space}
% ============================================================

\textbf{Why:} In the low-dimensional embedding we need a similarity measure that
(a) is symmetric by construction and (b) has \emph{heavier tails} than a
Gaussian. Heavy tails are crucial: they allow moderately dissimilar points to be
placed further apart without incurring a large penalty, alleviating the
``crowding problem'' (see Appendix~B).

t-SNE uses a \emph{Student's $t$-distribution with one degree of freedom}
(i.e.\ a Cauchy distribution) to define the joint similarity of the embeddings
$y_i, y_j$:
\[
  \boxed{q_{ij}
  \;=\;
  \frac{\left(1 + \|y_i - y_j\|^2\right)^{-1}}
       {\displaystyle\sum_{k \neq \ell}\left(1 + \|y_k - y_\ell\|^2\right)^{-1}}},
  \qquad q_{ii} = 0.
\]

This is symmetric by construction: $q_{ij} = q_{ji}$.

% ============================================================
\section*{Step 3: Formulate the Optimization Objective}
% ============================================================

\textbf{Why:} We want the distribution $Q = \{q_{ij}\}$ in embedding space to
match the distribution $P = \{p_{ij}\}$ in the original space as closely as
possible. The natural measure of how much one probability distribution diverges
from another is the \textbf{Kullback--Leibler (KL) divergence} (see Appendix~C).

The objective is:
\[
  \boxed{\mathcal{L}(Y)
  \;=\;
  \mathrm{KL}(P \,\|\, Q)
  \;=\;
  \sum_{i \neq j} p_{ij} \log \frac{p_{ij}}{q_{ij}}},
\]
where $Y = \{y_1, \dots, y_n\}$ collects all embedding coordinates.

Since $P$ is fixed (computed once from the data), minimising $\mathcal{L}$
over $Y$ is equivalent to maximising $\sum_{i\neq j} p_{ij}\log q_{ij}$, i.e.\
making $Q$ as large as possible wherever $P$ is large.

% ============================================================
\section*{Step 4: Derive the Gradient}
% ============================================================

\textbf{Why:} $\mathcal{L}$ has no closed-form minimiser; we use
\textbf{gradient descent} to iteratively update the embeddings. We therefore
need $\partial \mathcal{L} / \partial y_i$.

Let $Z = \sum_{k \neq \ell}(1 + \|y_k - y_\ell\|^2)^{-1}$ denote the
normalisation constant, so that
\[
  q_{ij} = \frac{w_{ij}}{Z}, \qquad w_{ij} = \left(1+\|y_i-y_j\|^2\right)^{-1}.
\]

Expand the loss, dropping terms constant in $Y$:
\[
  \mathcal{L}
  = \sum_{i\neq j} p_{ij}\log p_{ij}
  - \sum_{i\neq j} p_{ij}\log q_{ij}
  = \text{const} - \sum_{i\neq j} p_{ij}(\log w_{ij} - \log Z).
\]

\textbf{Differentiating with respect to $y_i$} (applying the chain rule):

\paragraph{Term 1: $\partial(\log w_{ij})/\partial y_i$.}

\[
  \frac{\partial \log w_{ij}}{\partial y_i}
  = \frac{1}{w_{ij}}\cdot\frac{\partial w_{ij}}{\partial y_i}
  = (1+\|y_i-y_j\|^2) \cdot (-1)(1+\|y_i-y_j\|^2)^{-2} \cdot 2(y_i-y_j)
  = -2\,w_{ij}(y_i-y_j).
\]

\paragraph{Term 2: $\partial \log Z / \partial y_i$.}

\[
  \frac{\partial \log Z}{\partial y_i}
  = \frac{1}{Z}\sum_{k\neq\ell}\frac{\partial w_{k\ell}}{\partial y_i}.
\]

Only terms with $k=i$ or $\ell=i$ are non-zero:
\[
  \frac{\partial \log Z}{\partial y_i}
  = \frac{1}{Z}\sum_{j\neq i}
    \left(\frac{\partial w_{ij}}{\partial y_i}+\frac{\partial w_{ji}}{\partial y_i}\right)
  = \frac{1}{Z}\sum_{j\neq i} 2 \cdot(-2\,w_{ij}(y_i-y_j))
  = -4\sum_{j\neq i}q_{ij}(y_i-y_j).
\]

(using $w_{ij}=w_{ji}$ and $q_{ij}=w_{ij}/Z$.)

\paragraph{Combining:}

\[
  \frac{\partial\mathcal{L}}{\partial y_i}
  = -\sum_{j\neq i}p_{ij}\left[\frac{\partial \log w_{ij}}{\partial y_i}
    - \frac{\partial \log Z}{\partial y_i}\right]
\]
\[
  = -\sum_{j\neq i}p_{ij}\Bigl[-2w_{ij}(y_i-y_j)\Bigr]
    + \sum_{j\neq i}p_{ij}\Bigl[-4\sum_{k\neq i}q_{ik}(y_i-y_k)\Bigr].
\]

Since $\sum_{j\neq i}p_{ij}=1$ (after symmetrisation and normalisation is
absorbed into $P$), the second term simplifies. Collecting and simplifying
(using the symmetry $p_{ij}=p_{ji}$, $q_{ij}=q_{ji}$):

\[
  \boxed{
  \frac{\partial \mathcal{L}}{\partial y_i}
  = 4\sum_{j\neq i}(p_{ij} - q_{ij})\,w_{ij}\,(y_i - y_j)
  }.
\]

\textbf{Interpretation:} Each pair $(i,j)$ contributes an attractive force
(when $p_{ij} > q_{ij}$, points are pulled together) or a repulsive force
(when $p_{ij} < q_{ij}$, points are pushed apart), weighted by the
heavy-tailed kernel $w_{ij}$.

% ============================================================
\section*{Step 5: Optimization via Gradient Descent}
% ============================================================

The embeddings are updated iteratively. Using \textbf{gradient descent with
momentum} (to accelerate convergence and reduce oscillations):

\[
  Y^{(t+1)}
  = Y^{(t)}
  - \eta\,\frac{\partial \mathcal{L}}{\partial Y}\bigg|_{Y^{(t)}}
  + \alpha(t)\Bigl(Y^{(t)} - Y^{(t-1)}\Bigr),
\]

where:
\begin{itemize}
  \item $\eta > 0$ is the \emph{learning rate},
  \item $\alpha(t) \in (0,1)$ is the \emph{momentum coefficient} at step $t$,
  \item $\partial\mathcal{L}/\partial Y$ stacks the per-point gradients from Step~4.
\end{itemize}

In practice, \textbf{early exaggeration} is applied during the first
$\sim\!250$ iterations: multiply all $p_{ij}$ by a factor (e.g.\ $12$) to
encourage the formation of tight, well-separated clusters before fine-tuning.

% ============================================================
\section*{Conclusion}
% ============================================================

t-SNE constructs a faithful low-dimensional embedding by:
\begin{enumerate}
  \item Converting high-dimensional pairwise distances to a symmetric probability
        distribution $P$ via per-point Gaussian kernels with perplexity-controlled
        bandwidths.
  \item Defining a low-dimensional similarity distribution $Q$ using a
        heavy-tailed Student-$t$ kernel to counter the crowding problem.
  \item Minimising the KL divergence $\mathrm{KL}(P\|Q)$ via gradient descent,
        where the gradient admits the elegant closed form
        \[
          \frac{\partial \mathcal{L}}{\partial y_i}
          = 4\sum_{j\neq i}(p_{ij} - q_{ij})\,w_{ij}\,(y_i - y_j).
        \]
\end{enumerate}

\begin{center}
\fbox{\parbox{0.88\textwidth}{\centering
t-SNE preserves local neighbourhood structure by matching a heavy-tailed
low-dimensional similarity distribution to a Gaussian high-dimensional one,
minimising their KL divergence through gradient descent.
}}
\end{center}

% ============================================================
\appendix
\section*{Appendix A: Perplexity and Bandwidth Selection}
% ============================================================

The bandwidth $\sigma_i$ for each point $x_i$ is chosen so that the conditional
distribution $P_{j\mid i}$ achieves a target \emph{perplexity}:
\[
  \mathrm{Perp}(P_{i})
  = 2^{H(P_{i})},
  \qquad
  H(P_{i}) = -\sum_{j\neq i} p_{j\mid i}\log_2 p_{j\mid i},
\]
where $H(P_i)$ is the Shannon entropy of the conditional distribution.
Perplexity can be interpreted as a smooth measure of the effective number of
neighbours. A typical value is $\mathrm{Perp} \in [5, 50]$. The bandwidth
$\sigma_i$ is found by binary search: increase $\sigma_i$ if the current
entropy is too low, decrease it if too high.

% ============================================================
\section*{Appendix B: The Crowding Problem}
% ============================================================

In high dimensions, the volume of a sphere of radius $r$ grows as $r^p$.
When projecting onto $d \ll p$ dimensions, there is simply not enough area to
faithfully represent all moderately distant neighbours---they get ``crowded''
into a small central region. Using a Gaussian in the low-dimensional space
would require placing moderately similar points so close together that they
overlap with genuinely similar ones.

The Student-$t$ kernel $w_{ij} = (1+\|y_i-y_j\|^2)^{-1}$ has a heavier tail
than a Gaussian: for large $\|y_i-y_j\|$, $w_{ij}$ decays as
$\|y_i-y_j\|^{-2}$ rather than exponentially. This allows moderate distances
in embedding space to correspond to moderate $q_{ij}$ values, relieving the
crowding.

% ============================================================
\section*{Appendix C: Kullback--Leibler Divergence}
% ============================================================

For two discrete probability distributions $P = \{p_{ij}\}$ and
$Q = \{q_{ij}\}$ over the same index set, the KL divergence is:
\[
  \mathrm{KL}(P \,\|\, Q)
  = \sum_{i \neq j} p_{ij} \log \frac{p_{ij}}{q_{ij}} \;\geq\; 0,
\]
with equality if and only if $P = Q$ (by Gibbs' inequality). KL divergence is
\emph{not} symmetric: $\mathrm{KL}(P\|Q) \neq \mathrm{KL}(Q\|P)$ in general.

In t-SNE, using $\mathrm{KL}(P\|Q)$ (rather than $\mathrm{KL}(Q\|P)$) has an
important asymmetric effect:
\begin{itemize}
  \item When $p_{ij}$ is large but $q_{ij}$ is small, the term
        $p_{ij}\log(p_{ij}/q_{ij})$ is very large $\to$ the loss strongly
        penalises placing nearby high-dimensional points far apart in the
        embedding.
  \item When $p_{ij}$ is small but $q_{ij}$ is large, the penalty is minimal
        $\to$ placing far-apart points nearby is less costly.
\end{itemize}
This asymmetry biases t-SNE towards faithfully representing local structure.

% ============================================================
\section*{Appendix D: Gradient Descent with Momentum}
% ============================================================

Standard gradient descent updates parameters as $\theta \leftarrow \theta - \eta
\nabla\mathcal{L}(\theta)$.
\textbf{Momentum} adds a fraction $\alpha$ of the previous update:
\[
  v^{(t+1)} = \alpha\,v^{(t)} - \eta\,\nabla\mathcal{L}(\theta^{(t)}),
  \qquad
  \theta^{(t+1)} = \theta^{(t)} + v^{(t+1)}.
\]
Momentum damps oscillations in directions of high curvature and accelerates
movement along consistent gradient directions, improving convergence speed and
stability.

\end{document}
